\begin{abstract}
Este trabajo aborda la clasificación automática de 104 tipos de flores mediante redes neuronales convolucionales (CNN) y técnicas de destilación de datos. Se utiliza el conjunto de imágenes de la competencia \textit{Petals to the Metal} de Kaggle, compuesto por 23,847 fotografías RGB distribuidas entre entrenamiento, validación y prueba. La similitud visual entre algunas clases, las variaciones de iluminación y el desbalance de los datos representan desafíos importantes para la clasificación.

Se implementó una CNN de cuatro bloques convolucionales (480,936 parámetros), entrenada mediante retropropagación con el optimizador Adam, parada temprana por F1 macro y guardado del mejor modelo. Además, se aplicó \textit{Distribution Matching} (DM) para generar conjuntos reducidos de imágenes sintéticas, y se comparó el entrenamiento con el conjunto completo, con subconjuntos aleatorios de imágenes reales y con conjuntos destilados de una y diez imágenes por clase.

Con el conjunto completo, el modelo alcanzó un F1 macro de 0.5538 y una exactitud de 61.61\,\% en validación. Las curvas de aprendizaje muestran que, a partir de la época 25, la pérdida de validación se estanca mientras la de entrenamiento sigue bajando, y que los conjuntos reducidos memorizan sus pocas imágenes. Con el conjunto completo se evaluaron tres variaciones de hiperparámetros (tasa de aprendizaje, \textit{dropout} y regularización): reducir la tasa de aprendizaje a $3\times10^{-4}$ produjo subajuste (F1 de 0.4605) y quitar la regularización aumentó el sobreajuste (F1 de 0.5332), mientras que la configuración de referencia obtuvo el mejor F1, 0.5784. Todas las tablas y figuras se reproducen con el notebook público del proyecto.
\end{abstract}

\begin{keywords}
Clasificación de imágenes \sep
Redes neuronales convolucionales \sep
Curvas de aprendizaje \sep
Destilación de datos \sep
F1 macro
\end{keywords}
